% =============================================================================
% Estado del Arte
% =============================================================================
\section{Estado del Arte}

% Análisis de los artículos consultados
\subsection{Análisis de Literatura}

El análisis de la literatura científica reciente muestra una evolución clara hacia la optimización de arquitecturas y la eficiencia operativa:

\begin{itemize}
\item \textcite{Shin2016} evaluaron el uso de arquitecturas como AlexNet y GoogLeNet en radiología torácica, demostrando que el aprendizaje por transferencia desde bases de datos generales (ImageNet) compensa eficazmente la escasez de imágenes médicas etiquetadas.

\item \textcite{MenchuLopez2025} realizó una revisión sistemática que sitúa a la radiología como la especialidad líder en adopción de IA, destacando que los algoritmos de deep learning ya igualan o superan la precisión de especialistas en tareas de cribado de cáncer de mama y gliomas.

\item \textcite{Laurent2025} documenta que a mediados de 2025 existen 873 algoritmos de IA autorizados por la FDA en radiología, con una velocidad de aprobación de aproximadamente 115 nuevos algoritmos por año. Según una encuesta de la ESR (2024) con 572 radiólogos europeos, el 48\% ya utiliza herramientas de IA (frente al 20\% en 2018), aunque solo el 2\% de las prácticas radiológicas estadounidenses han integrado plenamente lecturas asistidas por IA.

\item \textcite{PerezPelegri2023} profundizaron en la segmentación cardíaca, validando la arquitectura U-Net como el estándar para delimitar estructuras anatómicas y resaltando que el uso de modelos de conjunto (ensembles) es la mejor estrategia para evitar el sobreajuste (overfitting) en datos clínicos limitados.
\end{itemize}
% Comparación de enfoques metodológicos
\subsection{Comparación de Enfoques}

La siguiente tabla compara los enfoques de los estudios fundamentales analizados:

\begin{table}[ht]
\centering
\caption{Comparación de enfoques metodológicos en la literatura consultada}
\label{tab:comparacion}
\begin{tabularx}{\textwidth}{@{}lXXXX@{}}
\toprule
\textbf{Criterio} & \textbf{Shin et al. (2016)} & \textbf{Menchú López (2025)} & \textbf{Pérez-Pelegrí et al. (2023)} & \textbf{Laurent (2025)} \\
\midrule
\textbf{Modelo} & GoogLeNet / AlexNet (Transfer Learning) & CNN / IA Generativa / NLP & U-Net y Modelos de Conjunto (Ensembles) & Viz.ai, Mirai, Aidoc, GPT-4V \\
\textbf{Dataset} & ILD (Pulmón) y ganglios linfáticos & Revisión de 1,016 dispositivos autorizados & MRI Cardíaca (Eje corto) & 873 algoritmos FDA; encuesta ESR (572 radiólogos) \\
\textbf{Precisión} & 90.2\% de precisión en clasificación & AUC: 93.2\% (Gliomas) y 89.6\% (Mama) & Alto coeficiente Dice en segmentación & AUC $>$0.90 (Viz.ai); Sensibilidad $>$90\% (Aidoc) \\
\textbf{Limitaciones} & Escasez de datos médicos anotados & Sesgos algorítmicos y \enquote{AI drift} (deriva) & Sobreajuste en bases de datos pequeñas & Solo 2\% de adopción plena en EE.UU.; 80\% desconoce regulaciones \\
\bottomrule
\end{tabularx}
\end{table}

% Discusión de resultados reportados
\subsection{Discusión de Resultados}

Los resultados confirman que la IA basada en CNN es operativamente transformadora: las CNN identifican patologías críticas como el neumotórax en tiempo real, priorizando el triaje antes de que el radiólogo abra el estudio. No obstante, persiste una brecha importante entre el avance tecnológico y la preparación profesional: el 80\% de los radiólogos europeos desconoce las regulaciones vigentes para dispositivos de IA médica y ningún LLM ha sido aprobado por la FDA para uso clínico \parencite{Laurent2025}.
